\section{Cross-Encoders and Late Interaction}
\medskip

Bi-encoders trade expressiveness for efficiency. Cross-encoders do the opposite.

\bigskip
\subsection{Cross-Encoder}
\medskip

A cross-encoder feeds the pair jointly:
\[
\texttt{[CLS]} \; Q \; \texttt{[SEP]} \; D \; \texttt{[SEP]}.
\]

Full attention allows rich token-level interaction between query and document. However, at inference we must process every pair:
\[
O(N) \text{ forward passes per query.}
\]

This is computationally expensive, making cross-encoders suitable primarily for re-ranking a small candidate set.

\bigskip
\subsection{ColBERT and Late Interaction}
\medskip

ColBERT offers a middle ground. Queries and documents are encoded independently at the token level:
\[
\{T_i^Q\}, \quad \{T_j^D\}.
\]

Relevance is computed via
\[
\text{score}(Q, D)
=
\sum_i
\max_{j}
\langle T_i^Q, T_j^D \rangle.
\]

This preserves offline document encoding while enabling fine-grained token matching. The approach requires more memory than a single-vector bi-encoder, but avoids the full cost of cross-encoding. It balances scalability and expressiveness.


\bigskip

